% Planted-direction check -- Owner: Hassan, Red-teamer: Harsh
% Requires: \usepackage{booktabs}
% STATUS: provisional. Numbers are from 3 seeds. Circularity controls ARE run
% (disjoint-split + answer-absent); the literal out-of-library axis, which needs
% response generation + a GPU embedding pass, is not.
% Open items are marked [PROVISIONAL] in the text so they are easy to find.

\section{Planted-Direction Check}
\label{sec:planted}

The other analyses in this paper ask what the learned reward directions mean.
On real preference data there is no answer key, so a confident wrong label and a
correct one look the same. This section builds a case where the answer is known
in advance: we plant a concept, let the pipeline recover it, and check whether
the recovered direction can be named correctly.

The chain under test has four links: users who disagree along a known axis, an
mLoRe fit, a recovered reward direction, and an SAE readout that assigns a name.
A failure anywhere in that chain shows up as a wrong name, so the experiment is
designed to localise where it fails.

\subsection{Setup}

We construct synthetic users who disagree along six concept axes (\emph{fluency,
repetition, confidence, creativity, formatting, diversity}). For each axis, half
the users prefer high values of the concept and half prefer low, giving twelve
$(\text{concept},\text{sign})$ groups of twenty users each. Every user sees a
random $70\%$ subset of their group's preference pairs, so a single user's
direction carries subset noise on top of the concept; the group is therefore the
unit on which ground truth is defined, and we take one mean direction per group.
Each of the twelve directions now carries a label we know to be correct.

Preferences are expressed over real Skywork response embeddings, so the planted
signal is the only thing that is synthetic. The SAE is the frozen \texttt{D3}
checkpoint (dictionary size $16384$, $k=256$), trained on $41{,}730$ PRISM
response embeddings and never exposed to the planted axes.

We score the readout as forced choice: given a direction, the SAE must select
the correct $(\text{concept},\text{sign})$ pair from all candidates. Presence of
the right concept somewhere in the top features is not used, because the concept
library is internally correlated up to $|\cos| = 0.68$ and a permissive metric
would reward near-misses.

\subsection{Two controls}

\textbf{Null.} The same readout applied to random unit directions. This should
land at chance; if it does not, the scoring path is leaking information about
the answer rather than reading it off the direction.

\textbf{Ceiling.} The same readout applied to the \emph{true} concept vector
instead of the recovered one. This control is load-bearing here, because
training recovers which users are alike far better than it recovers what the
axes are: group-direction match is $0.989$, while subspace alignment is $0.755$
and best-axis match is $0.674$ (random floors $0.050$ and $0.030$). A recovered
direction is thus only a partial reconstruction of the concept it was planted
from. Without the ceiling, a weak readout could not be attributed to the SAE
rather than to the fit.

\subsection{Stage 1: can the SAE identify the planted concept?}

\begin{table}[h]
\centering
\begin{tabular}{lr}
\toprule
Direction supplied to the readout & Accuracy \\
\midrule
Recovered (planted concept)       & $0.778 \pm 0.079$ \\
True concept vector (ceiling)     & $0.667$ \\
Random unit directions (null)     & $0.083$ \\
\midrule
Chance ($1/12$)                   & $0.083$ \\
\bottomrule
\end{tabular}
\caption{Forced-choice identification over twelve $(\text{concept},\text{sign})$
groups. Mean over three seeds; per-seed accuracies $0.667$, $0.833$, $0.833$.}
\label{tab:planted-readout}
\end{table}

The readout identifies the planted concept at roughly nine times chance, and the
null sits exactly at chance, confirming the scoring path is not leaking. The
concept survives the full chain and the SAE can find it.

The ceiling falls slightly below the recovered condition. With twelve items per
seed the difference is one to two items and is within sampling noise; we do not
read it as the recovered direction being more informative than the vector it was
planted from. \emph{[PROVISIONAL: three seeds. To be re-run at ten seeds so this
comparison carries an interval rather than a verbal caveat.]}

\subsection{Stage 2: can the SAE name the concept from text?}

Stage 1 scores a direction against concept vectors, which are themselves
geometric objects. The task the method actually needs is harder: produce a name
from text. In Stage 2 the scoring path contains no concept vectors and no
embeddings. Feature meanings are derived from PRISM text alone, concept
signatures from contrast text alone, and the two are matched. We compare this
against a baseline that performs the same text matching \emph{without} routing
through SAE features.

\begin{table}[h]
\centering
\begin{tabular}{lr}
\toprule
Naming route & Accuracy \\
\midrule
Via SAE features       & $0.389$ \\
Without the SAE        & $0.417$ \\
\midrule
Null (SAE / no-SAE)    & $0.061$ / $0.066$ \\
Chance ($1/22$)        & $0.045$ \\
\bottomrule
\end{tabular}
\caption{Naming from text over twenty-two $(\text{concept},\text{sign})$ groups,
mean over three seeds. Per-seed: SAE $0.333, 0.500, 0.333$; no-SAE $0.500,
0.417, 0.333$.}
\label{tab:planted-naming}
\end{table}

Both routes beat chance by roughly eight times, so the planted concept is
recoverable from text. But the SAE route confers no measurable advantage over
the route that skips it. The two are also not independent: across the $36$
scored items they return the same $(\text{concept},\text{sign})$ prediction
$78\%$ of the time, and the discordant cells are small enough that the
comparison has almost no power to separate them --- one item correct under the
SAE route only, two under the baseline only, $13$ under both, $20$ under
neither. The defensible claim is therefore that routing through SAE features
gives \emph{no detectable advantage} over a simpler text baseline on this task,
not that it is worse.

\subsection{Is the concept library doing the work?}
\label{sec:planted-circularity}

In both stages the planted direction and the scoring map are built from the same
fifty contrast pairs per concept, so a hit could be self-matching rather than
concept reading. Two separable questions live inside that concern, and we test
each. A genuinely novel concept cannot be constructed from existing data --- it
requires generating new responses and embedding them --- so the second test
makes the planted concept \emph{unavailable} to the readout instead.

\paragraph{Disjoint text.} We split each concept's pairs into two halves, plant
users from half A only, and build the SAE concept profiles from half B only:
same concepts, no shared items. Half-A concept vectors agree with the full-set
vectors at $\cos = 0.992$, and user grouping is recovered perfectly
(group-direction match $1.000$), so the halves are individually adequate.

\begin{table}[h]
\centering
\begin{tabular}{lr}
\toprule
Planting set / scoring set & Accuracy \\
\midrule
Half A / half A (shared items)   & $0.694$ \\
Half A / half B (disjoint items) & $0.722$ \\
\midrule
Chance ($1/12$)                  & $0.083$ \\
\bottomrule
\end{tabular}
\caption{Identification accuracy with shared versus disjoint contrast items,
three seeds, twenty-five pairs per concept per half. Per-seed: shared $0.500,
0.833, 0.750$; disjoint $0.583, 0.750, 0.833$.}
\label{tab:planted-disjoint}
\end{table}

Accuracy does not depend on the planting and scoring sets sharing items --- the
disjoint condition is marginally higher, well within seed noise. Both sit below
the headline $0.778$ of Table~\ref{tab:planted-readout} because each half
carries half the data. The identification result in Stage~1 is therefore not
item memorisation.

\paragraph{Confidence when the answer is absent.} We then remove the planted
concept's own column from the candidate set, so the readout is forced to name a
direction whose true concept it does not have. From the readout's point of view
the planted axis is now outside its library. This is compared against removing a
\emph{random other} concept, leaving the correct answer available; both
conditions score over the same number of columns, so the confidence margin is
not inflated by the smaller denominator.

\begin{table}[h]
\centering
\begin{tabular}{lrr}
\toprule
Condition & Accuracy & Mean margin \\
\midrule
A random other concept removed (answer available) & $0.750$ & $0.406$ \\
The true concept removed (answer unavailable)     & ---     & $0.392$ \\
\midrule
\multicolumn{2}{l}{Ratio of margins, absent over present} & $0.966$ \\
\bottomrule
\end{tabular}
\caption{Readout confidence with and without the correct answer in the candidate
set, three seeds. Margin is the winning concept's share of total absolute
signature mass.}
\label{tab:planted-absent}
\end{table}

The readout is essentially as confident when the correct answer is unavailable
as when it is present: removing the true concept costs $3.4\%$ of the margin.
It does not decline to answer, and it does not signal that anything is wrong ---
it names a different concept with almost undiminished confidence.

This is the more consequential of the two controls. It means the margin carries
no information about whether a name is correct, and therefore that on real
directions --- where no answer key exists --- a confident SAE naming is not
evidence that the naming is right. Any downstream use of this readout needs an
external check on the label, not the readout's own confidence.

\emph{[PROVISIONAL: the literal pre-registered version --- planting an axis
generated outside the concept library and embedded through Skywork --- is not
run here. It requires response generation and a GPU embedding pass, and would
additionally test whether a novel axis is representable at all. The two controls
above address item-level circularity and answer-absence calibration, which are
the parts reachable from existing data.]}

\subsection{What this shows}

The SAE is not blind to preference directions. When a concept is planted and
scored geometrically, the SAE recovers it far above chance with a null that
behaves correctly (Table~\ref{tab:planted-readout}). But at the task the method
depends on --- turning a direction into a name using text --- passing through
SAE features buys nothing detectable over a simpler baseline
(Table~\ref{tab:planted-naming}). These are compatible: the information is
present and reachable, and the SAE decomposition is not what makes it readable.

The controls sharpen rather than soften this. Identification survives disjoint
text, so the Stage~1 result is real. But the readout is no less confident when
the correct answer has been removed from its candidate set
(Table~\ref{tab:planted-absent}), so it cannot flag its own failures. Taken
together: the SAE can find a planted concept, it adds nothing detectable to
naming one, and it gives no usable signal about when to believe it.

\subsection{Limitations}

The SAE operates off-distribution on the contrast set. Mean reconstruction error
there is $17.35$, against $11.98$ on held-out PRISM responses, so Stage~2 asks
the dictionary to describe inputs less like its training data than the
benchmark suggests.

Both stages are small. Stage~1 scores twelve items per seed and Stage~2 scores
twenty-two, over three seeds; intervals are correspondingly wide and the
SAE-versus-baseline contrast rests on three discordant items.

The answer-absent control removes a concept from the candidate set rather than
planting a genuinely novel one. A removed concept still correlates with what
remains in the library, so this is a lower bound on the confidence a truly
unrepresentable axis would attract, not a measurement of it.

Finally, the synthetic users disagree cleanly and by construction along a single
axis each. This is deliberate --- the point is an existence proof, and a
negative result under these conditions is more informative than a negative
result on noisy data --- but a positive result here is an upper bound on what to
expect from real preference data, not a prediction for it.

\appendix
\section{Reproduction details for the planted-direction check}
\label{app:planted}

Scripts are on branch \texttt{hassan/sae-planted-concept}. Stage~1 is
\texttt{PRISM/planted\_directions.py}, which builds users via
\texttt{PRISM/synthetic\_concept\_users.py} and emits one labelled group-mean
direction per $(\text{concept},\text{sign})$ group. Stage~2a is
\texttt{sae/scripts/planted\_concept\_readout.py} and Stage~2b is
\texttt{sae/scripts/planted\_concept\_naming.py}, the latter drawing feature
meanings from \texttt{sae/scripts/feature\_text\_profiles.py} over $3{,}381$
profiled features. The circularity controls of
Section~\ref{sec:planted-circularity} are
\texttt{sae/scripts/planted\_circularity\_controls.py}. Outputs are written to
\texttt{results/planted/}.

Two implementation points affect reproducibility.

First, differences are taken in \emph{activation} space: each response is
encoded separately and the codes subtracted, rather than encoding the difference
of the embeddings. The SAE encoder computes $\mathrm{enc}(x - b_{\text{pre}})$
with $b_{\text{pre}}$ the mean training embedding ($\|b_{\text{pre}}\| \approx
144$), so a difference vector, which is already centred near zero, would be read
as approximately $-b_{\text{pre}}$ and the top-$k$ selection would describe that
offset rather than the preference. The same caution applies to encoding any
unit-norm direction.

Second, all runs are pinned to CPU. \texttt{utils.py} resolves the device at
import and constructs parameters with \texttt{device=device}, so
\texttt{manual\_seed} draws a different initialisation stream under CUDA; on one
run this moved subspace alignment from $0.755$ to $0.795$ and best-axis match
from $0.674$ to $0.703$ with the data-side quantities unchanged. Figures in this
section are comparable across machines only when CUDA is disabled.
